
\newpage

\section{Interfacciamento PC--Board: API Python e Protocollo Seriale}


\subsection{Architettura del Sistema di Sincronizzazione}

Il firmware non ha connettività di rete diretta. La sincronizzazione dei dati meteo e il recupero della posizione
GPS avvengono tramite la porta seriale USB, che il Pico espone come CDC
(Communication Device Class).

Il flusso è di tipo \textbf{master-slave}: il PC chiede, la board risponde o riceve; non c'è polling autonomo dalla board verso il PC.

\begin{figure}[h]
	\centering
	\begin{tabular}{lll}
		\textbf{PC} & & \textbf{RP2040} \\
		\hline
		\texttt{pcb\_bridge.py} & $\rightarrow{\texttt{GET\_GPS\textbackslash n}}$ & \texttt{sys\_manager\_handle\_serial()} \\
		& $\leftarrow{\texttt{GPS\_DATA:LIVE,lat,lon}}$ & \\
		Nominatim API & & \\
		OpenMeteo API & & \\
		& $\rightarrow{\texttt{WXC,City,T,W,H,C\textbackslash n}}$ & \\
		& $\leftarrow{\texttt{OK\_WXC}}$ & \\
	\end{tabular}
	\caption{Sequenza di scambio dati tra \texttt{pcb\_bridge.py} e firmware.}
\end{figure}

\subsubsection*{Protocollo di comunicazione}

Il protocollo creato é in puro ASCII, ed ogni messaggio è terminato da \texttt{'\textbackslash n'}.

%% MATTIA da verificare l'ultimo
\begin{tabularx}{\textwidth}{@{}l l X@{}}
	\toprule
	\textbf{Direzione} & \textbf{Messaggio} & \textbf{Significato} \\
	\midrule
	PC $\to$ Board & \texttt{GET\_GPS} & Richiede posizione corrente \\
	Board $\to$ PC & \texttt{GPS\_DATA:LIVE,lat,lon} & Fix GPS valido, coordinate correnti \\
	Board $\to$ PC & \texttt{GPS\_DATA:CACHE,lat,lon} & Nessun fix, coordinate da Flash \\
	PC $\to$ Board & \texttt{WXC,City,T,W,H,C} & Pacchetto meteo (5 campi CSV) \\
	Board $\to$ PC & \texttt{OK\_WXC} & Conferma ricezione meteo \\
	PC $\to$ Board & \texttt{GET\_STATUS} & Richiesta stato diagnostico \\
	Board $\to$ PC & (multiriga) & Dump errori e stato sistema \\
	\bottomrule
\end{tabularx}


\subsection{Lo Script Python}

\texttt{tools/api-weather/pcb\_bridge.py} è il programma lato PC che
orchestra l'intera sequenza. Le uniche dipendenze esterne sono \texttt{pyserial} e \texttt{requests} per le API HTTP.

\begin{lstlisting}[style=py, caption={pcb\_bridge.py --- flusso principale}]
	def main():
	ser = serial.Serial("/dev/ttyACM0", 115200, timeout=5)
	time.sleep(5)          # attende il boot del firmware post-reset USB
	ser.reset_input_buffer()
	
	# 1. Richiede la posizione GPS alla board
	ser.write(b"GET_GPS\n")
	
	data_line = None
	for _ in range(50):    # attende fino a 50 righe in risposta
	line = ser.readline().decode('utf-8', errors='ignore').strip()
	if line.startswith("GPS_DATA:"):
	data_line = line.replace("GPS_DATA:", "")
	break
	
	# data_line e' "LIVE,lat,lon" oppure "CACHE,lat,lon"
	source, lat, lon = data_line.split(',')[0], *map(float, data_line.split(',')[1:])
	
	# 2. coordinate -> nome citta'
	city = get_city_name(lat, lon)   # Nominatim (OpenStreetMap)
	
	# 3. Meteo dalla posizione
	cur = fetch_weather(lat, lon)    # OpenMeteo API
	
	# 4. Formatta e invia alla board
	wxc = f"WXC,{city},{cur['temperature_2m']},{cur['wind_speed_10m']},"
	wxc += f"{cur['relative_humidity_2m']},{cur['weather_code']}\n"
	ser.write(wxc.encode('utf-8'))
\end{lstlisting}

\textbf{NOTA 1}: \texttt{time.sleep(5)} dopo l'apertura della porta seriale é necessario poiché
quando il PC apre la connessione CDC, il Pico riceve un segnale DTR che interpreta come richiesta di reset. Il firmware riparte da capo, esegue l'intera sequenza di boot (LittleFS, display, animazione) e solo dopo
qualche secondo il sistema è in uno stato stabile pronto a processare comandi.
Senza questo delay, \texttt{GET\_GPS} verrebbe inviato durante il boot e
ignorato o perso nel buffer.


\textbf{NOTA 2}: Il loop \texttt{for \_ in range(50)} legge fino a 50 righe cercando quella
che inizia con \texttt{GPS\_DATA:}. Infatti le righe precedenti possono contenere
output di debug del firmware che vengono silenziosamente scartate. 

\subsubsection*{API HTTP esterne}

\begin{lstlisting}[style=py, caption={pcb\_bridge.py --- geocoding e meteo}]
	def get_city_name(lat, lon):
	# Nominatim: reverse geocoding con OpenStreetMap
	url = "https://nominatim.openstreetmap.org/reverse"
	params = {"lat": lat, "lon": lon, "format": "json", "zoom": 10}
	headers = {"User-Agent": "PPSE-Lab-Weather-Bridge"}
	r = requests.get(url, params=params, headers=headers, timeout=5)
	address = r.json().get("address", {})
	return address.get("city") or address.get("town") or address.get("village") or "Unknown"
	
	
	# OpenMeteo
	def fetch_weather(lat, lon):
	url = "https://api.open-meteo.com/v1/forecast"
	params = {
		"latitude": lat, "longitude": lon,
		"current": "temperature_2m,relative_humidity_2m,"
		"wind_speed_10m,weather_code",
	}
	r = requests.get(url, params=params, timeout=10)
	return r.json()["current"]
\end{lstlisting}


\subsubsection*{Formattazione del pacchetto WXC}

La stringa inviata al firmware ha il formato:
\begin{center}
	\texttt{WXC,\textit{City},\textit{temp\_C},\textit{wind\_kmh},\textit{humidity\_pct},\textit{wmo\_code}}
\end{center}


Lato firmware, il parsing in \texttt{sys\_manager\_handle\_serial()} usa
\texttt{indexOf} su Arduino \texttt{String} per trovare i cinque separatori:

\begin{lstlisting}[language=C++, caption={system\_manager.cpp --- parsing WXC}]
	int p1 = cmd.indexOf(',');       // dopo "WXC"
	int p2 = cmd.indexOf(',', p1+1); // dopo City
	int p3 = cmd.indexOf(',', p2+1); // dopo temp
	int p4 = cmd.indexOf(',', p3+1); // dopo wind
	int p5 = cmd.indexOf(',', p4+1); // dopo humidity
	
	if (p5 != -1) {  // pacchetto completo
		cmd.substring(p1+1, p2).toCharArray(w.city, 32);
		w.temp_ext   = cmd.substring(p2+1, p3).toFloat();
		w.wind_speed = cmd.substring(p3+1, p4).toFloat();
		w.humidity   = cmd.substring(p4+1, p5).toInt();
		w.weather_code = cmd.substring(p5+1).toInt();
		w.valid = true;
		sys_manager_update_weather(w);
		Serial.println("OK_WXC");
	}
\end{lstlisting}

Se uno qualsiasi dei cinque separatori manca (\texttt{p5 == -1}), il pacchetto
viene scartato senza risposta. Lo script Python interpreta l'assenza di
\texttt{OK\_WXC} entro il timeout come errore.
 
